\section*{Overview}

Offline queries are the answer when online order is a burden rather than a requirement. If all queries are known in
advance, I can sort or bucket them into a friendlier order and maintain a data structure state that only moves forward.

\section*{When to Use It}

Use offline processing when:

\begin{itemize}
  \item all queries are available before answering begins,
  \item sorting queries by one parameter makes updates monotone,
  \item the online version would require a more complex data structure.
\end{itemize}

\section*{Core Idea}

Choose an order in which processing becomes incremental. A common example is:

\begin{itemize}
  \item sort array elements by value,
  \item sort range queries by threshold \(x\),
  \item add all elements with value \(\le x\) to a Fenwick tree before answering that query.
\end{itemize}

\section*{Key Insight}

Offline means I am allowed to change time order as long as I remember the original query indices. That freedom often
replaces a hard dynamic problem with a simple sweep.

\section*{Operations / Main Technique}

\begin{itemize}
  \item annotate each query with its original index,
  \item sort by a parameter that should move monotonically,
  \item sweep through updates and answer queries in that order,
  \item restore the original answer order at the end.
\end{itemize}

\paragraph{Worked example.}
To count how many values in \([l, r]\) are at most \(x\), sort the queries by \(x\), insert positions of array values
in increasing-value order into a Fenwick tree, and query the active-count range when each threshold is reached.

\section*{Correctness Intuition}

The sorted sweep ensures that, before answering a query with threshold \(x\), the data structure exactly represents all
contributions that should be active for that threshold and none that should not. The original query order is irrelevant
once the answers are stored by index.

\section*{Complexity Analysis}

Typical offline sweeps cost \(O((n + q)\log n)\) after sorting, instead of paying for a heavier online structure.

\section*{Implementation}

The sample code implements the threshold-query pattern described above with a Fenwick tree. It answers range counts of
values \(\le x\).

\section*{Common Pitfalls}

\begin{itemize}
  \item Forgetting to preserve original query indices.
  \item Sorting by the wrong key so the sweep state no longer matches the query meaning.
  \item Forcing an offline solution when the statement requires interactive or online answers.
\end{itemize}

\section*{Variants / Extensions}

\begin{itemize}
  \item Mo's algorithm for reordering by block locality.
  \item Parallel binary search.
  \item CDQ divide and conquer for time-ordered offline events.
\end{itemize}

\section*{Practice Problems}

\begin{itemize}
  \item Count values \(\le x\) in subarrays.
  \item Offline rectangle counting after coordinate compression.
  \item Process add-and-query events after sorting by time or threshold.
\end{itemize}

\section*{References}

\begin{itemize}
  \item \href{https://codeforces.com/catalog}{Codeforces Catalog}.
  \item \href{https://cp-algorithms.com/data_structures/fenwick.html}{cp-algorithms: Fenwick Tree} for the standard sweep helper.
\end{itemize}
